\documentclass[11pt,a4paper,twoside]{article}
\usepackage[utf8]{inputenc}
\usepackage[english]{babel}
\usepackage{actuarialangle, actuarialsymbol}


\begin{document}

Define a discount function and state some immediate properties


A discount function is a positive valued function v, of two nonnegative
variables, satisfying 
\( v(s, t)v(t, u) = v(s, u) \)
for all values of \( s, t, u \).
Some immediate Corollaries from the Definition are:

\( v(s, s) = 1 \) so
\( v(s, t)v(t, s) = v(s, s) = 1 \)
so
\( v(s, t) = v(t, s)^{-1} \)
In particular for discrete time units, the recursive relation 
\( v(n) = v(n-1)v(n-1, n), v(0) = 1\), should hold.
Although we have called \( v \) a ‘discount’ function, the common English usage of the word
really applies to the case where \( s < t \). In that case, \( v(s, t) \) will be normally less than \( 1 \), 
and the function is returning the discounted amount of \( 1 \) unit paid at a later date.
Some like therefore to refer to the case \( t < s \) as accumulation function.
Typically \( v \) is expressed with respect to one variable, that is \( v(t) = v(0, t) \),
which can be bootstraped from the original values via 
\( v(s, t) = \frac{v(t)}{v(s)} \) and the recursion formula.
In the course it was mostly assumed that \( v(s, t) = v^{t-s} \), for all \( t \).


Define the related components 
involved in future/present value calculations
to the discount function

The rate of interest for the time interval \( k \) to \( k + 1 \) is the quantity
\( i_k = v(k+1, k) - 1 \)
The rate of discount for the time interval \( k \) to \( k + 1 \) is the quantity
\( d_k = 1 - v(k, k+1) \)

which are related to the discount function via 

\( d_k = i_kv(k, k+1) = \frac{i_k}{1 + i_k} \)
\( i_k = d_kv(k+1, k) = \frac{d_k}{1 - d_k} \)
\( v(k+1, k) = \frac{1}{1 + i_k} \)
\( \frac{1}{d_k} - \frac{1}{i_k} = 1 \)

If the unit of time is continuous one speaks instead of the force of interest,
given by
\( \delta_t = \frac{v'(t)}{v(t)} = \frac{d}{dt}\ln{v(t)} \),
with the assumption that \( v(t) = e^{\int_0^t \delta_\tau d\tau}\)
under a constant force of interest this simplifies to
\( e^{\delta} = 1 + i \), with \( \delta = \ln{1 + i} \)

In case that on is interested in the intra-annually/nominal rates, one has
\(\left( 1 + \frac{i^{(m)}}{m} \right)^m = 1 + i \left( 1 - \frac{d^{(m)}}{m} \right)^m = 1 - d\)
\(i^{(m)} = m \left[ (1 + i)^{\frac{1}{m}} - 1 \right] d^{(m)} = m \left[ 1 - (1 - d)^{\frac{1}{m}} \right]\)
and the modified discount factor is
\( (\frac{1}{(1 + i)})^{\frac{1}{m}} = v^{\frac{1}{m}} \)


State the future/present values of different forms of annuities

The future value was also called terminal value in the lecture

The basic formula's are \\
PV of $ 1 $ Due in $ t $ Years: \\
$ PV = \frac{1}{a(t)} = (1 + i)^{-t} = v^t = e^{-\delta t} = (1 - d)^t $ \\
$ = \left( 1 + \frac{i^{(m)}}{m} \right)^{-mt} = \left( 1 - \frac{d^{(m)}}{m} \right)^{mt} $ \\
$ FV = a(t) = (1 + i)^t = e^{\delta t} = (1 - d)^{-t} = \left( 1 + \frac{i^{(m)}}{m} \right)^{mt}$ \\
$\left( 1 - \frac{d^{(m)}}{m} \right)^{-mt}$ \\

Annuity-Immediate (in advance) \\
(PV one period before first payment, FV at time of last payment) \\

$ PV = \ax{\angl n} = \sum_{i=1}^n v^i = \frac{1 - v^n}{i} $ \\
$ FV = \sx{\angl n} = \sum_{j=0}^{n-1} (1 + i)^j = \frac{(1 + i)^n - 1}{i} = \ax{\angl n}(1 + i)^n$ \\
Annuity-Due (in arrears) \\
(PV at time of first payment, FV one period after last payment) \\
$PV = \ax**{\angl n} = \sum_{j=0}^{n-1} v^j = \frac{1-v^n}{d} = (1 + i)\ax{\angl n}$ \\
$ = (\frac{i}{d}\ax{\angl n} = 1 + a{\angl{n -1}}$ \\

$ FV = \sx**{\angl n} = \sum_{j=0}^{n} (1 + i)^j = \frac{(1 + i)^n - 1}{d}$ \\
$ = (1+i)\sx{\angl n} = (\frac{i}{d})\sx{\angln} = \sx{\angl{n+1}} - 1 $ \\

and the intra monthly versions are

in advance \\
$ \ax**{\angl n}[(m)] = \sum^{nm-1}_{t = 0} \frac{1}{m}v^{\frac{t}{m}} = \frac{1}{m}\frac{(1 - v^n)}{1 - v^{\frac{1}{m}}}$ \\
in arrears \\
$\ax{\angl n}[(m)] = \sum^{nm}_{t = 1} \frac{1}{m}v^{\frac{t}{m}} = \frac{1}{m}v^{\frac{1}{m}}\frac{(1 - v^n)}{1 - v^{\frac{1}{m}}}$ \\

In case of a deferred annuity \\

Annuity-Immediate (in advance) \\
$ \ax**[m|]{\angl n} = \sum_{t=m}^{n + m -1} v^t = v^m\frac{1 - v^n}{1 - v}$ \\

Annuity-Due (in arrears) \\
$\ax[m|]{\angl n} = \sum_{t=m+1}^{n + m} v^t = v^m \frac{1 - v^n}{i}$ \\
$= \ax**[m+1|]{\angl n} = v^m \ax{\angl n} = \ax{\angl{m+n}} - a{\angl{m}}$ \\



State the expressions for perpetuities

in advance\\
$ \ax**{\angl \infty} = \sum_{t=0}^\infty v^t = \lim_{n \to \infty} \ax**{\angl n} = \frac{1}{1 - v} = \frac{1}{d}$
in arrears\\
$ \ax{\angl \infty} = \sum_{t=1}^\infty v^t = v\ax**{\angl \infty} = \frac{1}{i}$




Give a few real-life examples to calculate annuities

A pensioner wants to receive for the next $20$ years a yearly annuity,
paid in advance, with $R = 12 000$ What is the capital
requirement $A$ if $i = 0.035$? \\
$ A = R\ax**{\angl{20}}$

A married couple won $100 000$ Euros in the lottery. They decide to
buy a yearly immediate annuity in arrears for the next $10$ years.
How high is $R$, if the yearly interest rate is $i = 0.04$? \\
$R = \frac{10^5}{\ax{\angl{10}}}$


Answer in a nutshell why insurance

Insurance is an agreement to transfer costs associated with an uncertain event.
- Rachel Z. Friedman
that is exchanging something uncertain (loss, damage, litigation) with sth. certain (premium payments).
And it collectivises risks (law of large numbers)
Mutual and social insurance assume that if enough individuals combine together, paying into
the common pool an amount determined by their risks, they can equitably share the burdens
of a misfortune that happens to strike any one of them. The fairness of this agreement hinges
on the parties’ mathematical equality.



Define the discrete time surplus process
The (discrete-time) surplus process of the insurer with initial capital \(c_0 \ge 0\) is defined as
\(C_t = C_t^{c_0} := c_0 + \sum_{u = 1}^t (\pi_u - S_u), t \in \mathbb{N}_0 \left( \sum_{u = 1}^0 \cdot \cdot \cdot := 0 \right)\).

Here, \(\{(\pi_u, S_u)|u \in \mathbb{N}\}\) is a family of random variables with the following properties:

\(\{(\pi_u, S_u)|u \in \mathbb{N}\}\) independent
\((\pi_u, S_u) \stackrel{d}{=} (\pi_1, S_1), u \in \mathbb{N}\)
\(\mathbb{P}(S_1 \ge 0) = 1\)
\(\mathbb{P}(\pi_1 > 0) = 1\) and there is \(r > 0\) such that \(M_{\pi_1}(s) := \mathbb{E}[e^{s \pi_1}] < \infty s \in (-r, r)\)

The last equation depends on a moment-generating function for \( \pi_1 \) which has positive radius of convergence.
and is equivalent to \( \mathbb{P}(\pi_1 > 0) \) and \( \exists r > 0: \mathbb{E}[e^{r\pi_1}] < \infty \).


Some remarks
\( C_t \) models equity/net asset value at time \( t \in \mathbb{N}_0 \)
As long as \( C_t \geq 0 \) the insurer is solvent; otherwise, if \( C_t < 0 \), the insurer cannot cover its
liabilities \( \implies \) bankruptcy/ruin!

Thus \( \mathbb{P}(C_t < 0) \) should be kept as small as possible.

\( \mathbb{P}\circ(\pi_t - S_t)^{-1} \) (distribution of \( \pi_t - S_t \) under \( \mathbb{P} \)) 
independent of \( t \in \mathbb{N}_0 \), which means time homogeneity.
The argument comes form \( \mathbb{P} \circ (\pi_t, S_t)^{-1} \) not depending on 
\( t \) and \( \mathbb{R}^2 \ni (x, y) \mapsto x - y \) is measurable.

The family \( (\pi_t - S_t),\ t \in \mathbb{N}_0 \) is i.i.d.
\( (C_t)_{t \in \mathbb{N}_0} \) is a random walk starting in \( c_0 \geq 0 \).



Give the definition of ruin time

Let \(C = (C_t)_{t \in \mathbb{N}_0}\) be the surplus process.

\(\tau: \Omega \rightarrow \mathbb{N}_0 \cup \{\infty\},\)
\(w \mapsto \inf \{s \in \mathbb{N}_0 | C_s(w) < 0\} \)
is the ruin time of \(C\), with (\inf \theta = \infty).

Remark
Classical concept from probability theory: stopping times
\(\tau\) is stopping time w.r.t. filtration
\(\mathbb{F} := (\sigma(C_u|0 \le u \le t))_{t \in \mathbb{N}_0}\) 
canonical filtration of \(C\))

In particular: \(\tau\) is an \(\mathbb{F}\)-mb. random variable
The associated probabilities are

Ruin probability up to time \(t \in \mathbb{N}\):
\(\psi_t(c_0) := \mathbb{P}(\tau \le t) = \mathbb{P}(\exists s \in \{0, 1, ..., t\}: C_s < 0)\)

Ultimate ruin probability:
\(\psi(c_0) = \mathbb{P}(\tau < \infty)\)

Some intuitive propositions then are
Let \( c_0 \geq 0 \)
1. \( 0 ≤ s ≤ t \implies 0 \leq \psi_s(c_0) \leq \psi_t(c_0) \leq 1\)
2. \( \psi(c_0) = \mathbb{P}(\exists t \in \mathbb{N}_0) : C_t^{c_0} < 0) = \lim_{t \to \infty} \psi_t(c_0)\).
and for
\( 0 \leq c_0 \leq \tilde{c}_0 \), \( \psi_t(\tilde{c}_0) \leq \psi_t(c_0), t \in \mathbb{N}_0, \psi(\tilde{c}_0) \leq \psi(c_0) \).


State the Net Profit Condition (NPC) and the motivating Proposition.


The Proposition is:
\(S_1\) integrable (\(\pi_1\) integrable by definition)
\(\mathbb{P}(S_1 \ne \pi_1) > 0\)

Then the following assertions hold:
\(\mathbb{E}[S_1] < \mathbb{E}[\pi_1] \implies \lim_{t \rightarrow \infty} C_t = \infty\) \(\mathbb{P}\)-a.s. 
due to (law of large numbers)
\(\mathbb{E}[S_1] > \mathbb{E}[\pi_1] \implies \lim_{t \rightarrow \infty} C_t = -\infty\) \(\mathbb{P}\)-a.s. 
due to (law of large numbers)
\(\mathbb{E}[S_1] = \mathbb{E}[\pi_1] \implies \mathbb{P}(\lim \sup_{t \rightarrow \infty} C_t = -\lim \inf_{t \rightarrow \infty} C_t = \infty) = 1\) 
(Borel-Cantelli Lemma)

from which the NPC follows, namely
\( \mathbb{E}[\pi_1] > \mathbb{E}[S_1] \),
note that this does not guarantee no bankruptcy, it just ensures the absence of certain bankruptcy.


State the recursive formula's for ruin probabilities

Let \(c_0 \ge 0\).
\(\forall t \in \mathbb{N}_0: \psi_{t + 1}(c_0) = 1 - F(c_0) + \int_{-\infty}^{c_0} \psi_t(c_0 - y) dF(y)\)
\(\psi(c_0) = 1 - F(c_0) + \int_{-\infty}^{c_0} \psi(c_0 - y)dF(y)\)
where \( F \) is the cumulant distribution function of \( S_1 - \pi_1 \).


Define Lundberg coefficients and state their key properties

Suppose \( \exists R > 0 \), s.t. \( \mathbb{E}[e^{S_1 - \pi_1)R}] = 1 \),
then \( R \) is called the Lundberg coefficient.
Lemma:
If Lundberg coefficient exists, then it is unique.
Theorem:
Suppose the NPC holds and there exists a Lundberg coefficient, then \( \psi(c_0) \leq e^{-Rc_0}, c_0 \geq 0 \).


Motivate what makes a Lundberg coefficients powerful in the context of insurance.

\(\psi(c_0) \le e^{-R c_0} \rightarrow 0, c_0 \rightarrow \infty\), exponentially fast!
\(\forall \epsilon > 0 \exists c_0 > 0\) s.t. \(\psi(c_0) \le \epsilon\) "too large"
Don't have to choose \(c_0\)
Ultimate ruin probability made arbitrarily small by choosing sufficient initial capital

If \(S_j\) and \(\pi_1\) ind. (e.g. \(\pi_1 \in \mathbb{R}_+\)), then \(S_1 \wedge k\) and \(\pi_1\) ind., \(k \in \mathbb{N}\)
Assume Lundberg coefficient \(R\) exists

\(\implies \mathbb{E}[e^{R(S_1 \wedge k)}] \mathbb{E}[e^{-R \pi_1}] = \mathbb{E}[e^{R(S_1 \wedge k - \pi_1)}] \le \mathbb{E}[e^{R(S_1 - \pi_1)}] = 1\)

Monotone conv. \(\implies \mathbb{E}[e^{R S_1}] = \lim_{k \rightarrow \infty} \mathbb{E}[e^{R(S_1 \wedge k)}] \le \mathbb{E}[e^{-R \pi_1}]^{-1} < \infty\)
\(\implies\) claim size distribution \(\mathbb{P}_{S_1}\) is light-tailed, i.e., survival function \(\mathbb{P}(S_1 > x)\) decays exponentially fast in \(x\):
\(\mathbb{P}(S_1 > x) \le \text{const.} \cdot e^{-Rx}\)
Markov ineq.

However this is a problematic assumption about many distributions in insurance applications.


Define subexponentiality and state how it relates the Lundberg bounds

\(H: \mathbb{R} \rightarrow [0, 1]\) cdf such that \(H(0) = 0\). \(H\) is subexponential if
\[\lim_{x \rightarrow \infty} \frac{1 - H^{*2}(x)}{1 - H(x)} = 2.\]
\(H^{*2}\) cdf of self-convolution of \(H\), i.e., cdf of \(Z = X + Y\) for two i.i.d. random variables \(X, Y\) with \(F_X = F_Y = H\)

An easier condition to verify this is regular variation, namely a function \( f : \mathbb{R}_{+} \to \mathbb{R}_{+} \)
is regularly varying if for all \( a \in \mathbb{R}_{+} \) it holds
\( \lim{x \to \infty} \frac{f(ax)}{f(x)} = a^\gamma \), for some \( \gamma \geq 0 \).

Suppose \( H \) is subexponential. Then
\( \forall r > 0 : \lim_{x \to \infty} e^{rx}(1 - H(x)) = \infty \), 
then one can deduce with Choquets representation of the expectation
\( \mathbb{E}[X] = \int_{0}^\infty \mathbb{P}(X > x)dx \),
that
\( \mathbb{E}[e^{r(S_1 - \pi_1)}] = \mathbb{E}[e^{r\pi_1} r\int_{\mathbb{R}} e^{rk}(1 - H(k)) dk = \infty\),
where \( k = \frac{\log(x)}{r}, dk = \frac{dx}{rx} \),
thus Lundberg coefficients can't exists.

Famous examples that caught by this are 
- many Weibull distributions
- log-normal distributions
- log-gamma distributions
- Pareto distributions

which do not satisfy this

\end{document}
